The linear representation hypothesis --- that high-level concepts are encoded as directions in the residual stream of large language models --- underpins core interpretability techniques such as linear probing, steering vectors, and activation editing.
Yet whether these directions \emph{compose} reliably via vector addition remains poorly understood.
We systematically test the compositionality of 25 concept categories from \bats spanning morphological, encyclopedic, and lexicographic relations, extracting directions from the unembedding matrix of \pythia and evaluating them with both geometric metrics (causal inner products, a new Composition Fidelity Score) and functional steering experiments.
Our central finding is that compositionality is not a property of direction \emph{pairs} but of direction \emph{quality}: morphological concepts form strong linear directions (leave-one-out consistency $0.38 \pm 0.12$) that compose freely, encyclopedic concepts are moderate ($0.35 \pm 0.09$), and lexicographic relations such as synonymy and antonymy mostly fail to form coherent linear directions at all (many LOO $< 0.1$).
Among well-formed directions, composition via vector addition is remarkably robust --- functional steering with composed directions $\vd_A + \vd_B$ achieves $>$96\% accuracy across all tested pair types, with minimal cross-concept interference.
These results draw a clear empirical boundary around where additive composition of linear directions can and cannot be trusted.
